\pdfoutput=1
\documentclass[11pt]{article}

\usepackage{acl}
\usepackage{times}
\usepackage{latexsym}
\usepackage[T1]{fontenc}
\usepackage[utf8]{inputenc}
\usepackage{microtype}
\usepackage{inconsolata}
\usepackage{graphicx}
\usepackage{booktabs}
\usepackage{amsmath}
\usepackage{amssymb}
\usepackage{amsfonts}
\usepackage{algorithm}
\usepackage{algorithmic}
\usepackage{enumitem}
\usepackage{multirow}
\usepackage{float}

\title{Adaptive Computation Control for Efficient Reasoning \\ in Large Language Models}

\author{
  Neill Rawani \\ 2025202027 \And
  Tanisha Upadhyay \\ 2025201088 \And
  Aryan \\ 2025201017
}

\begin{document}
\maketitle

\begin{abstract}
Large Language Models (LLMs) achieve state-of-the-art performance on complex mathematical and logical benchmarks, largely driven by step-by-step "Chain-of-Thought" (CoT) prompting. However, standard prompting techniques inherently lead to "over-computation." Models predictably continue generating extensive reasoning chains even after the correct answer has been resolved internally, driven by learned conversational distributions rather than logical necessity. This redundancy dramatically increases token usage, exacerbates memory bandwidth bottlenecks during autoregressive decoding, and limits real-world deployment scalability. This paper proposes a model-agnostic Adaptive Reasoning Control System that dynamically determines optimal stopping points during generation. Crucially, our approach requires no modifications to the base model's pre-trained weights and does not rely on internal, often miscalibrated, confidence logits. We formulate the reasoning process as a sequential decision-making problem, utilizing a lightweight Multi-Layer Perceptron (MLP) trained on dense semantic embeddings of intermediate reasoning traces. Our automated data-labeling pipeline captures empirical reasoning distributions, allowing the controller to synchronously monitor the generative state. Upon breaching a calibrated confidence threshold, the controller executes a structural verification protocol and forces a hard exit. We demonstrate empirically that our synchronous neural controller preserves high logical accuracy (92.0\% on GSM8K, 76.0\% on MathQA) while reducing redundant token generation by up to 17.51\%.
\end{abstract}

\section{Introduction and Motivation}
The advent of Large Language Models (LLMs), such as \texttt{Qwen2.5-Math-7B-Instruct}, has revolutionized automated reasoning. By employing "Chain-of-Thought" (CoT) prompting \cite{wei2022chainofthought}, these models can decompose complex queries into intermediate logical steps, effectively trading increased computational time for higher accuracy. However, this profound capability incurs a steep, often unjustified, computational cost. 

During inference, autoregressive token generation is heavily memory-bandwidth bound. Maintaining the Key-Value (KV) cache for extensive, multi-step reasoning traces consumes significant VRAM, and the strictly sequential nature of token generation inherently limits batch parallelization. 

The core motivation for this project is grounded in the principle of dynamic computational allocation: reasoning complexity is heterogenous, and therefore, computational budgets should not be static. A simple arithmetic word problem might be unequivocally solved by the third reasoning step. Yet, constrained by autoregressive pre-training objectives that prioritize conversational completeness and politeness, the LLM will natively ramble up to 60 steps, generating hundreds of unnecessary tokens purely to satisfy structural norms or execute redundant "double-check" routines.

Our objective is to design, implement, and evaluate an "Adaptive Computation Strategy" via a Learning-Based Early-Exit Controller. Operating synchronously during the decoding process, our lightweight controller analyzes the semantic trajectory of the running text trace. If the controller determines that the correct answer is already present within the context window, it triggers an "Early Exit." This strategy aims to decouple reasoning accuracy from token volume, proving that highly capable black-box LLMs can be externally regulated, saving massive amounts of compute and GPU time.

\subsection{Summary of Contributions}
In this work, we make the following core contributions to the field of adaptive inference:
\begin{itemize}[leftmargin=*]
    \item We introduce a novel, model-agnostic generation-level early-exit framework that treats the foundational reasoning model entirely as a black box, requiring no parameter fine-tuning.
    \item We design an automated, self-labeling data generation pipeline (Phase A) capable of translating unconstrained reasoning traces into mathematically bounded, discrete verification steps, featuring robust trace filtration algorithms.
    \item We demonstrate empirically across multiple benchmarks that our synchronous neural controller preserves high logical accuracy (92.0\% on GSM8K, 76.0\% on MathQA) while reducing redundant token generation by up to 17.51\%.
    \item We provide a rigorous qualitative error analysis, formalizing the failure modes of regex-based verification within early-exit frameworks.
\end{itemize}

\section{Background and Related Work}
Adaptive computation in neural networks is an expanding domain, aiming to optimize the Pareto frontier of accuracy and inference latency. Our proposed system builds upon, and explicitly contrasts with, several recent paradigms in the field.

\subsection{Mathematical Formulation of CoT}
Standard autoregressive language modeling seeks to maximize the probability of a token sequence $Y = (y_1, y_2, \dots, y_T)$ given an input prompt $X$. The probability is factored as:
\begin{equation}
    P(Y | X) = \prod_{t=1}^{T} P(y_t | X, y_{<t})
\end{equation}
In Chain-of-Thought (CoT) reasoning \cite{wei2022chainofthought}, the sequence $Y$ is artificially prolonged into intermediate steps $S$ before reaching the final answer $A$. The generation terminates only when the model naturally emits an \texttt{EOS} token at step $T$. Our system formally intervenes to truncate this product at $t = \tau_{exit} \ll T$, forcibly injecting an \texttt{EOS} state when the semantic representation of $y_{<\tau_{exit}}$ implies $A$ is fully resolved.


\subsection{The Bottleneck of Autoregressive Decoding}
To understand the necessity of token-level early exits, one must formalize the computational cost. For a transformer layer processing a sequence of length $s$ with a batch size of $b$, the self-attention mechanism requires caching the Key and Value tensors. The memory footprint scales linearly:
\begin{equation}
    M_{KV} \approx 2 \times b \times s \times L \times h \times d_{head} \times P
\end{equation}
As $s$ grows during a protracted CoT trace, memory bandwidth becomes the primary bottleneck. The time to decode a single token is fundamentally constrained by the time required to load the KV cache from High-Bandwidth Memory (HBM) to the streaming multiprocessors. Because $M_{KV}$ grows with every generated token, late-stage token generation is significantly slower than early-stage generation. Therefore, forcefully reducing the generated sequence length $\Delta s$ provides a highly disproportionate, super-linear improvement to total system throughput.

\subsection{Architectural Layer Exits vs. Generation Exits}
Historically, adaptive inference focused on internal structural modifications. Early-exit frameworks like DeeBERT \cite{xin2020deebert} append auxiliary classifiers to intermediate transformer layers, halting forward passes if internal hidden states reach sufficient confidence. While effective for standard classification tasks, these methods require modifications to pre-trained weights and fail to address autoregressive token over-generation across time steps. Our approach shifts the paradigm from internal layer-depth exits to temporal external generation exits.

\subsection{Heuristic vs. Learned Control Markers}
To achieve generation-level dynamic exits, contemporary methods often rely on monitoring specific linguistic artifacts or heuristic markers. For instance, DEER \cite{yang2025dynamic} monitors the occurrences of explicit "Wait" tokens to gauge internal reasoning states. Concurrently, EAT \cite{wang2025entropy} tracks next-token entropy variance after specific structural tags (e.g., \verb|</think>|) to identify when the model transitions from active reasoning to a finalized output. These heuristic approaches are inherently brittle, highly sensitive to prompt phrasing, and require base models specifically fine-tuned to emit these explicit markers. In contrast, our Learning-Based Controller evaluates the holistic semantic representation of the reasoning trace using a dense retrieval model, immunizing the system against static syntactic variations.

\subsection{Budgeting Signals and Ex-Ante Allocation}
Recent literature supports predicting computational needs \textit{ex-ante} (prior to generation) to bypass the KV bottleneck entirely. \citet{snell2025learning} proposed an input-adaptive allocation framework that predicts the inherent difficulty of a query to dynamically assign an optimal computing budget. Similarly, \citet{xie2025steering} demonstrated methodologies to steer LLM thinking through explicit budget guidance. These findings theoretically validate our planned \textbf{Difficulty-Aware Static Budgeting} baseline (Phase 3). However, we hypothesize that purely predictive ex-ante approaches suffer from high variance compared to our synchronous, state-aware MLP controller.

\subsection{Consistency and Convergence-Based Halting}
Output stability serves as a robust theoretical signal for stopping. Traditional Self-Consistency \cite{wang2022selfconsistency} improves Chain-of-Thought reasoning by sampling $K$ independent paths, though this aggressively multiplies computational complexity to $\mathcal{O}(K \times N)$. Recent works by \citet{chen2025answer} and \citet{liu2025early} demonstrate that answer convergence—when a model's internal conclusion stops altering across sequential steps—often precedes the formal conclusion of the textual chain. This literature directly validates our planned \textbf{Consistency-Based Early Exit} baseline (Phase 4), which seeks to approximate this convergence detection organically along a single $\mathcal{O}(N)$ generation path without the overhead of multi-path sampling.

\section{Dataset Details and Prompt Engineering}
We leverage distinct mathematical reasoning benchmarks representing varying distributions of problem complexity. A summary of the dataset domains and empirical baselines is provided in Table \ref{tab:dataset_stats}.

\begin{table}[h]
\centering
\small
\begin{tabular}{@{}llcc@{}}
\toprule
\textbf{Dataset} & \textbf{Domain} & \textbf{Target Format} & \textbf{Base Tokens} \\ \midrule
GSM8K & Grade Math & Exact Integer & 386.0 \\
MathQA & Adv. Logic & Alphabetical & 500.0 \\
SVAMP & Perturbations & Exact Integer & 600.0* \\ \bottomrule
\end{tabular}
\caption{Dataset characteristics and unconstrained baseline token expenditures. (*SVAMP baseline is an estimated fallback upper bound).}
\label{tab:dataset_stats}
\end{table}

\subsection{GSM8K Benchmark}
GSM8K comprises high-quality grade school math word problems requiring 2 to 8 steps of basic arithmetic. Due to its lower inherent complexity, GSM8K is a prime candidate for early exiting. To enforce explicit mathematical bounding, we utilize the following few-shot template:

\begin{quote}
\footnotesize
\textbf{System Prompt (GSM8K):} \\
Solve the following math problem step-by-step. State your final numerical answer clearly. After stating the answer, briefly double-check your work to ensure it is correct. \\
\textbf{Example Problem:} Natalia sold clips to 48 of her friends in April, and then she sold half as many clips in May. How many clips did Natalia sell altogether in April and May? \\
\textbf{Example Solution:} Natalia sold 48 clips in April. In May, she sold half as many, which is 48 / 2 = 24 clips. Altogether, she sold 48 + 24 = 72 clips. The final answer is 72. Double-check: 48 (April) + 24 (May) = 72. Correct.
\normalsize
\end{quote}

\subsection{MathQA Benchmark}
MathQA represents a large-scale dataset of advanced mathematical queries structured strictly as multiple-choice outputs. The reasoning paths here are inherently protracted, elevating the empirical baseline token expenditure to 500.0 tokens. We strictly prompt the model to output alphabetical answers:

\begin{quote}
\footnotesize
\textbf{System Prompt (MathQA):} \\
Solve the following multiple-choice math problem step-by-step. State your final numerical answer clearly, and explicitly mention the correct option letter. After stating the answer, briefly double-check your work to ensure it is correct. \\
\textbf{Example Problem:} a shopkeeper sells an article at a loss of 10 \% . if he had sold it for rs . 45 more , he would have gained 5 \% . find the cost price of the article ? \\
Options: a ) 250 , b ) 300 , c ) 350 , d ) 400 , e ) 450 \\
\textbf{Example Solution:} Let the cost price be x. Loss = 10\% of x = 0.1x. Selling price = x - 0.1x = 0.9x. If sold for 45 more, new selling price = 0.9x + 45. New gain = 5\% of x = 0.05x. New selling price = x + 0.05x = 1.05x. Therefore, 0.9x + 45 = 1.05x. 0.15x = 45 => x = 45 / 0.15 = 300. The cost price is 300. The correct option is b.
\normalsize
\end{quote}

\section{Proposed Methodology and Architecture}
Our architecture enforces adaptive generation by treating the foundational \texttt{Qwen2.5-Math-7B-Instruct} as an un-modified black box. 

\subsection{Phase A: Automated Trace Generation}
\label{sec:phase_a}
The objective of Phase A is to construct an empirical dataset of sequential reasoning steps mapped to binary correctness targets. The base model autoregressively generates a sequence of tokens $Y$, chunked into discrete steps. 

At every step $S_i$, a deterministic verification oracle scans the cumulative text bounds for the known true target answer, $A_{true}$. Because LLMs format answers unpredictably, this oracle utilizes a cascaded regex algorithm. For exact numerical matches, we deploy negative lookbehind and lookahead assertions to prevent substring collision:
\begin{equation}
    \text{Regex Pattern: } \texttt{(?<!\textbackslash d)} A_{true} \texttt{(?!\textbackslash d)}
\end{equation}

Furthermore, specifically for the GSM8K extraction logic, we implement a delimiter-aware fallback. If the model natively generates the \texttt{"\#\#\#\#"} boundary typical of the dataset's ground truth, the oracle strictly isolates the trailing numerical string:
\begin{algorithm}[H]
\caption{GSM8K Ground Truth Oracle Extraction}
\label{alg:gsm_extract}
\begin{algorithmic}[1]
\REQUIRE Generation $S_i$
\IF{\texttt{"\#\#\#\#"} $\in S_i$}
    \STATE $\text{tail} \leftarrow \text{split}(S_i, \texttt{"\#\#\#\#"})[-1]$
    \STATE $nums \leftarrow \text{regex\_findall}(\texttt{r'-?\textbackslash d+(?:\textbackslash .\textbackslash d+)?'}, \text{tail})$
    \IF{$nums$}
        \STATE \textbf{return} $nums[0]$
    \ELSE
        \STATE \textbf{return} $\text{tail}$
    \ENDIF
\ENDIF
\STATE \textbf{return} \texttt{""}
\end{algorithmic}
\end{algorithm}

The automated generation and labeling logic assigns binary ground truth states to the sequence:
\begin{algorithm}[H]
\caption{Phase A: Automated Trace Generation}
\label{alg:trace_gen}
\begin{algorithmic}[1]
\REQUIRE Dataset $\mathcal{D}$, Max Steps $N=60$, Model $\mathcal{M}$
\STATE Initialize Empty Corpus $\mathcal{C}$
\FOR{each $(X, A_{true})$ in $\mathcal{D}$}
    \STATE $S_0 \leftarrow X$, $\text{Solved\_Flag} \leftarrow \text{False}$
    \STATE Trace $\leftarrow []$
    \FOR{$i = 1$ \TO $N$}
        \STATE $S_i \leftarrow \mathcal{M}.\text{generate\_step}(S_{i-1}, \text{limit=20})$
        \IF{\textbf{not} $\text{Solved\_Flag}$}
            \IF{$f_{verify}(S_i, A_{true})$}
                \STATE $\text{Solved\_Flag} \leftarrow \text{True}$
                \STATE $y_i \leftarrow 1$
            \ELSE
                \STATE $y_i \leftarrow 0$
            \ENDIF
        \ELSE
            \STATE $y_i \leftarrow 1$
        \ENDIF
        \STATE Trace.append( ($S_i, y_i$) )
        \IF{$\mathcal{M}.\text{emitted\_eos}()$}
            \STATE \textbf{break}
        \ENDIF
    \ENDFOR
    \IF{$\text{Solved\_Flag}$}
        \STATE $\mathcal{C}.\text{append}(\text{Trace})$
    \ENDIF
\ENDFOR
\STATE \textbf{return} $\mathcal{C}$
\end{algorithmic}
\end{algorithm}

\subsection{Data Filtration and Ground Truth Labelling}
To ensure the neural controller learns true logical resolution rather than coincidental string matching, we implement a rigorous trace filtration protocol. A trace is discarded entirely if the answer was generated suspiciously early. Specifically, if \texttt{solved\_at\_step} $\le 2$, the trace is dropped, as this statistically indicates the model hallucinated the correct numerical string prior to executing substantive mathematical logic:
\begin{equation}
    \mathcal{D}_{filtered} = \{ Y \in \mathcal{D} \mid y_{final} = 1 \land \text{Solved\_Step} > 2 \}
\end{equation}

\subsection{Phase B: Controller Training \& Architecture}
With the filtered ground-truth dataset established, Phase B transitions to training the learning-based controller.

\paragraph{Semantic Embedding Extraction:} We map the raw text of each step $S_i$ into a continuous vector space utilizing a frozen \texttt{sentence-transformers/all-MiniLM-L6-v2} encoder, denoted as $E(\cdot)$. The dense semantic embedding is extracted as $v_i \in \mathbb{R}^{384}$. The \texttt{MiniLM-L6} architecture guarantees that embedding extraction time remains negligible compared to the autoregressive token generation time.

To provide spatial and temporal context, we concatenate two scalar values: the current step index $i$, and the local token count for that specific step segment, $t_i$. 
\begin{equation}
    z_i = [v_i \oplus i \oplus t_i] \in \mathbb{R}^{386}
\end{equation}
To ensure gradient stability, the scalar subset of $z_i$ is standard-scaled using a pre-fitted normalizer ($\mu, \sigma$) to produce $\hat{z}_i$.

\paragraph{MLP Controller Architecture:} The early-exit mechanism is parameterized by a Multi-Layer Perceptron projecting from $\mathbb{R}^{386} \rightarrow \mathbb{R}^1$. To prevent internal covariate shift across the highly correlated step-wise features, we implement Batch Normalization (\texttt{BatchNorm1d}) and rigorous Dropout:
\begin{align}
    h_1 &= \text{Dropout}_{0.3}(\text{ReLU}(\text{BN}_{128}(W_1 \hat{z}_i + b_1))) \\
    h_2 &= \text{Dropout}_{0.3}(\text{ReLU}(\text{BN}_{64}(W_2 h_1 + b_2))) \\
    p_i &= \sigma(W_3 h_2 + b_3)
\end{align}
The output $p_i \in (0,1)$ represents the confidence that the correct answer is fully generated. 

\begin{table}[h!]
\centering
\small
\begin{tabular}{@{}ll@{}}
\toprule
\textbf{Hyperparameter} & \textbf{Value} \\ \midrule
Input Dimension & 386 (384 Sem + 2 Scalar) \\
Hidden Layers & 128, 64 \\
Activation Functions & ReLU, Sigmoid (Output) \\
Dropout Probability & 0.3 \\
Batch Normalization & Applied pre-activation \\
Optimizer & AdamW \\
Learning Rate & 0.001 \\
Weight Decay & 1e-5 \\
Loss Function & Binary Cross Entropy (BCE) \\
Batch Size & 256 \\
Epochs & 20 (with early stopping) \\ \bottomrule
\end{tabular}
\caption{Hyperparameter configuration for the MLP Controller.}
\label{tab:hyperparams}
\end{table}

\paragraph{Offline Controller Training Dynamics:}
During Phase B execution, the filtered dataset was randomized and split into an 80\% training set and a 20\% validation test set. A critical challenge in training the step-wise controller is class imbalance; because the correct answer only appears at the tail end of a reasoning trace, the vast majority of steps $S_i$ are correctly labeled as $0$ (False). 

To ensure the model did not simply collapse into majority-class prediction, we rigorously logged Precision and Recall at a binarization threshold of $0.5$:
\begin{align}
    \text{Precision} &= \frac{TP}{TP + FP} \\
    \text{Recall} &= \frac{TP}{TP + FN}
\end{align}
The BCE loss naturally penalizes premature positive predictions, teaching the MLP to require high semantic certainty before outputting $p_i \to 1$, thus mitigating False Positives that would lead to catastrophic premature exits during live inference.

\paragraph{Answer Verification Protocol:} During live evaluation, if $p_i \geq \tau$ (e.g., $0.85$), the system invokes a rigorous verification layer to prevent truncation before the LLM has formatted its logic. The cascaded \texttt{has\_clear\_answer} function verifies the presence of structural bounding (e.g., \verb|\boxed{X}|) or explicit semantic declarations. If formatting is absent, the controller suppresses the halt command.

\begin{algorithm}[H]
\caption{Adaptive Early Exit Generation Loop}
\label{alg:early_exit}
\begin{algorithmic}[1]
\REQUIRE Prompt $X$, Max Steps $N=60$, Threshold $\tau$, Model $\mathcal{M}$, Controller $\mathcal{C}$, Extractor $E$
\STATE $S_0 \leftarrow X$
\FOR{$i = 1$ \TO $N$}
    \STATE $S_i, t_i \leftarrow \mathcal{M}.\text{generate\_step}(S_{i-1}, \text{limit=20})$
    \STATE $v_i \leftarrow E.\text{encode}(S_i)$
    \STATE $\hat{z}_i \leftarrow \text{normalize}(v_i \oplus i \oplus t_i)$
    \STATE $p_i \leftarrow \mathcal{C}.\text{forward}(\hat{z}_i)$
    \IF{$p_i \geq \tau$}
        \IF{\text{has\_clear\_answer}($S_i$)}
            \STATE \textbf{return} \text{extract\_final\_answer}($S_i$)
        \ELSE
            \STATE \textbf{continue} \COMMENT{Force formatting}
        \ENDIF
    \ENDIF
\ENDFOR
\STATE \textbf{return} \text{extract\_final\_answer}($S_N$) 
\end{algorithmic}
\end{algorithm}

\section{Proposed Evaluation Metrics}
\paragraph{Token Savings Ratio (TSR):} The aggregate percentage of tokens saved by the adaptive controller relative to the empirical baseline length ($B_{dataset}$) across the test corpus $\mathcal{D}_{test}$:
\begin{equation}
    \text{TSR} = \frac{1}{|\mathcal{D}_{test}|} \sum_{j=1}^{|\mathcal{D}_{test}|} \left( 1 - \frac{T_{adaptive}^{(j)}}{B_{dataset}} \right) \times 100\%
\end{equation}

\paragraph{Relative Accuracy Retained (RAR):} To verify that compute savings do not catastrophically degrade logic:
\begin{equation}
    \text{RAR} = \frac{\text{Accuracy}_{adaptive}}{\text{Accuracy}_{baseline}} \times 100\%
\end{equation}

\section{Empirical Results and Evaluation (Phase 5)}
With the multi-layer perceptron controller trained on a corpus of unconstrained reasoning traces, the pipeline underwent empirical evaluation to map the degradation curve of RAR against TSR.

\begin{table}[h!]
\centering
\small
\begin{tabular}{@{}lcccc@{}}
\toprule
\textbf{Dataset} & \textbf{Acc. (\%)} & \textbf{Base Tokens} & \textbf{Adaptive} & \textbf{TSR (\%)} \\ \midrule
GSM8K & 92.00 & 386.0 & 318.4 & \textbf{17.51} \\
MathQA & 76.00 & 500.0 & 426.0 & \textbf{14.80} \\ \bottomrule
\end{tabular}
\caption{Empirical evaluation results operating at a confidence threshold of $\tau = 0.85$.}
\label{tab:results_summary}
\end{table}

\subsection{GSM8K Benchmark Results}
Our evaluation on a \texttt{test} split of 25 GSM8K problems yielded highly promising results. Operating synchronously with a confidence threshold of $\tau = 0.85$, the pipeline achieved an absolute accuracy of \textbf{92.0\%} while registering an aggregate Token Savings Ratio (TSR) of \textbf{17.51\%}. 

\begin{figure}[t]
    \centering
    \includegraphics[width=\linewidth]{token_expenditure_comparison.pdf}
    \caption{Comparison of actual token expenditure using the adaptive controller versus the empirical baseline average (386 tokens) across the 25 GSM8K test samples. Samples are sorted by adaptive token usage.}
    \label{fig:token_comparison}
\end{figure}

These findings confirm our primary hypothesis: highly capable autoregressive models generate substantial conversational redundancy on simpler arithmetic domains. By continuously evaluating the semantic state of the context window, our controller successfully truncated generation on the majority of queries without compromising logical correctness.

\begin{figure}[t]
    \centering
    \includegraphics[width=\linewidth]{token_savings_distribution.pdf}
    \caption{Distribution of absolute token savings per query on the GSM8K test split. The mean token savings demonstrates a significant leftward shift in computational budget.}
    \label{fig:token_dist}
\end{figure}

\subsection{MathQA Benchmark Results}
To evaluate robustness on advanced logical constraints, we evaluated the pipeline on the MathQA dataset, which necessitates strict multiple-choice alphabetical outputs. The unconstrained baseline ($B_{math\_qa}$) was established at 500.0 tokens per query.

Our evaluation yielded an absolute accuracy of \textbf{76.00\%} with a TSR of \textbf{14.80\%}. The average tokens generated before exit was recorded at 426.0 tokens. While absolute accuracy naturally degrades on the more difficult MathQA set, the sustained $\sim$15\% token savings illustrates that the controller remains highly effective at mitigating redundancy even in protracted environments. 

\section{Qualitative Error Analysis \& Case Studies}
To better understand the behavioral dynamics across multiple formatting domains, we conducted a rigorous qualitative analysis of specific inference traces.

\subsection{Successful Early Exits (Arithmetic Completion)}
The system demonstrated high proficiency in recognizing completed arithmetic logic, bypassing the LLM's propensity to re-verify its own work. Consider the following trace excerpt from a GSM8K query:

\begin{quote}
\small
\textit{"...Janet sells each remaining egg for \$2. Therefore, her total revenue per day is 9 $\times$ 2 = 18 dollars. The final answer is \textbackslash boxed\{18\}. Double-check: 16 (total eggs) - 7 (eggs eaten and baked) = 9 (eggs left), and 9"}
\normalsize
\end{quote}

In unconstrained settings, the model natively continues to generate the remainder of the "Double-check" paragraph. Our controller recognized the semantic finality, triggered a hard exit at exactly 300 tokens, and successfully saved 86.0 tokens without altering final accuracy.

\subsection{Successful Early Exits (Multiple-Choice)}
The controller successfully adapted to the multiple-choice formatting specific to MathQA. In a complex problem involving the calculation of pastry costs, the LLM generated:

\begin{quote}
\small
\textit{"...Therefore, the total cost for all the pastries is 204 + 160 + 330 = 694 dollars. Since 694 corresponds to option c, the final answer is option c."}
\normalsize
\end{quote}

Despite lacking explicit LaTeX markers, the controller correctly identified the semantic resolution of the logic and verified the alphabetical formatting constraints (\verb|\b(?:option|answer)\s+([a-e])\b|). It executed an early exit at 331 tokens, saving 169.0 tokens (a 33.8\% reduction).

\subsection{Extended Runways (Dynamic Allocation)}
A critical observation in the distribution is the presence of "negative savings". Rather than a failure, this highlights the advantage of dynamic allocation. In a complex GSM8K query regarding travel times, the LLM required an extended chain of logic to isolate the variables:

\begin{quote}
\small
\textit{"...Total distance back = 0 miles + 15 miles + 120 miles = 135 miles. Finally, let's determine how far John is from home at the end of the 4 hours: Distance from home = Distance away - Total distance back = 180 miles - 135 miles = 45 miles. The final answer is \textbackslash boxed\{45\}."}
\normalsize
\end{quote}

The LLM required 480 tokens to fully resolve this logic. A Static Budgeting baseline would have prematurely truncated this query at exactly 386 tokens, guaranteeing a failed resolution. The neural controller correctly identified the unresolved semantic state and dynamically granted an extended computational runway.

\subsection{Premature Truncation (False Positives)}
The primary vector for accuracy degradation stems from false-positive exits induced by formatting ambiguities. In one instance, the model generated the correct intermediate logic, concluding with:
\begin{quote}
\small
\textit{"Total earnings = 400 dollars + 60 dollars = 460 dollars. The final answer is \textbackslash boxed\{46"}
\normalsize
\end{quote}
The controller triggered a hard exit at 360 tokens, extracting the incorrect partial answer "46" instead of the true answer "460". The structural verification protocol recognized the beginning of the LaTeX bound (\verb|\boxed{|) but failed to verify its closure (\verb|\}|), demonstrating a critical need for stricter structural regex requirements.

\begin{algorithm}[H]
\caption{Cascaded Extraction Protocol ($f_{verify}$)}
\label{alg:extraction}
\begin{algorithmic}[1]
\REQUIRE Text trace $S$, True Answer $A_{true}$
\IF{Regex Match \texttt{\textbackslash\textbackslash boxed\textbackslash\{([\textasciicircum\}]+)\textbackslash\}} in $S$}
    \STATE $E \leftarrow \text{Extract bounding contents}$
    \STATE \textbf{return} CheckBoundedNumerical($E$, $A_{true}$)
\ENDIF
\IF{Substring \texttt{"the final answer is"} in $S$}
    \STATE $E \leftarrow \text{Extract trailing string}$
    \STATE \textbf{return} CheckTrailingNumerical($E$, $A_{true}$)
\ENDIF
\IF{Substring \texttt{"double-check:"} in $S$}
    \STATE $E \leftarrow \text{Extract preceding string}$
    \STATE \textbf{return} CheckPrecedingNumerical($E$, $A_{true}$)
\ENDIF
\STATE \textbf{return} NaiveNumericalExtraction($S$, $A_{true}$)
\end{algorithmic}
\end{algorithm}

\section{Future Work and Scope}
While our initial methodology establishes a robust foundation for generation-level early exiting, several avenues remain for optimizing and practically expanding the Adaptive Reasoning Control System. 

\subsection{Performance Optimization: Accuracy and Compute}
A primary focus of our future work is to structurally increase both the Relative Accuracy Retained (RAR) and Token Savings Ratio (TSR). Currently, the system operates on a fixed static confidence threshold ($\tau = 0.85$). We plan to implement dynamic, instance-aware thresholds that adjust based on initial semantic entropy and mathematical complexity. By extensively fine-tuning the MLP controller's hyperparameters and utilizing higher-dimensional embedding extraction, we aim to mathematically minimize false-positive premature exits. This optimization will allow us to push RAR significantly closer to 99\% while continuously accelerating inference, potentially scaling TSR beyond 25\%. 

\subsection{Cross-Dataset and Cross-Model Generalization}
Our current evaluation protocols are constrained to strict mathematical reasoning paradigms (GSM8K and MathQA) using the \texttt{Qwen2.5-Math-7B-Instruct} base model. Future project scope will evaluate the controller's generalization capabilities across wider linguistic, Boolean, and multi-hop logical domains.
\begin{itemize}[leftmargin=*]
    \item \textbf{Additional Evaluation Datasets:} We will extend our training and testing pipeline to broader question-answering benchmarks, specifically \texttt{StrategyQA} (for logical boolean induction) and \texttt{HotpotQA} (for multi-hop semantic reasoning). Unlike math datasets, which guarantee explicit final bounding delimiters, these linguistic datasets will require entirely different oracle mechanics to grade.
    \item \textbf{Diverse Base Models:} A pressing research question is whether early-exit dynamics trained exclusively on Qwen traces can zero-shot generalize or transfer effectively to other open-weights mathematical architectures (e.g., Llama-3-Instruct-8B, Mistral-7B). We will rigorously benchmark the MLP's ability to evaluate semantic reasoning tokens across disjoint models without necessitating Phase B retraining.
\end{itemize}

\subsection{Phase 3: Difficulty-Aware Static Budgeting}
The first alternative architectural modification we will explore is an ex-ante static budgeting mechanism. A fundamental limitation of the synchronous dynamic controller is the step-by-step latency overlay of computing dense semantic representations at runtime. To bypass this, a lightweight regressor $f_{static}(\cdot)$ will be trained to predict the required computational budget \textit{before} the LLM begins generation, utilizing only the features of the input prompt heuristic $\phi(X)$ (e.g., syntactic tree depth or raw word count):
\begin{equation}
    \hat{N} = \lfloor f_{static}(\phi(X)) \rceil
\end{equation}
The base LLM will subsequently be strictly constrained to generate exactly $\hat{N}$ tokens without any dynamic interruptions. While we theorize this purely predictive approach will suffer from notably higher variance and potential task collapse, it will act as an essential benchmark for our dynamic methodology.

\subsection{Phase 4: Consistency-Based Early Exit}
Our second alternative future approach will explore mathematical stability signaling without requiring any external mapping networks. We will implement a rolling consistency checker configured with a stability window of dimension $m$. Under this framework, if the internally extracted answer entity remains structurally and mathematically identical across $m$ consecutive autoregressive reasoning steps:
\begin{equation}
    A_i \equiv A_{i-1} \equiv \dots \equiv A_{i-m+1}
\end{equation}
the generation loop is automatically terminated without consulting the MLP. This approach organically harnesses stability as a proxy for internal confidence, bypassing the requirement for offline embedding generation and producing a robust, purely heuristic-driven counterpart to our learning-based architecture.

\section{Limitations}
While the Adaptive Reasoning Control System presents significant efficiency gains, it is subject to several methodological limitations. Primarily, the Answer Verification Protocol relies heavily on predictable prompt adherence. If the base LLM dynamically alters its output format, the controller defaults to conservative over-generation, mitigating savings. Furthermore, synchronously executing a semantic embedding and an MLP forward pass at every reasoning step introduces a minor latency overhead that marginally offsets the clock-time gained from token reduction.

\section{Broader Impact and Ethical Considerations}
The deployment of Large Language Models carries an immense environmental cost due to the electricity required for prolonged GPU inference. Techniques that natively reduce the token expenditure required to arrive at a logical conclusion—without necessitating massive, energy-intensive model fine-tuning—contribute directly to "Green AI" initiatives. By actively truncating redundant autoregressive generation, our adaptive control system mathematically reduces the carbon footprint per user query, facilitating more sustainable and equitable access to advanced reasoning capabilities in resource-constrained environments.

\section{Conclusion}
In this paper, we presented a model-agnostic Adaptive Reasoning Control System designed to optimize the computational efficiency of Large Language Models. By framing reasoning as a step-wise sequence and deploying a synchronous neural controller trained on empirical trace embeddings, we successfully decoupled logical accuracy from autoregressive token exhaustion. Our initial evaluations demonstrated that a black-box LLM can maintain 92.0\% accuracy on GSM8K and 76.0\% on MathQA while achieving up to a 17.51\% reduction in total token generation. This paradigm shift from static conversational padding to state-aware dynamic allocation represents a crucial step toward scalable, latency-optimized deployment of advanced reasoning models.

\bibliography{custom}
\bibliographystyle{acl_natbib}

\end{document}
